\section{Карта на верификационните сценарии и репозиторийните артефакти}
\label{sec:verification_appendix}

Настоящото приложение систематизира връзката между архитектурните твърдения в дипломната работа и конкретните артефакти в хранилището. Целта му не е да дублира вече изложените аргументи, а да предостави компактна и проверима карта на това къде в кода и тестовата инфраструктура се намират основните доказателства за коректност.

Подобно приложение е особено полезно в контекста на настоящата разработка, защото библиотеката не е просто концептуален prototype. Тя е реално изградена като CMake проект със source код, unit тестове, integration тестове и conditional live тестове за различни backend-и. Именно тази материална структура позволява на дипломната работа да премине отвъд ниво ``архитектурно намерение'' и да бъде защитена като инженерно верифицируема система.

\subsection{Обща структура на доказателствения слой}

Верификацията в хранилището следва архитектурната структура на библиотеката. Production кодът е концентриран главно под \path{lib/include/ORM/}, а тестовете са разделени между локален unit слой и integration/live слой. CMake файловете регистрират тези тестове по различен начин в зависимост от това дали даден backend е винаги наличен, условно наличен или изисква външна среда.

\begin{figure}[H]
\centering
\begin{tikzpicture}[node distance=1.35cm and 1.25cm]
    \node[ormaccent] (src) {\path{lib/include/ORM/}\\production code};
    \node[ormblock, below left=of src] (unit) {\path{tests/unit/}\\локални invariants};
    \node[ormblock, below right=of src] (integration) {\path{tests/integration/}\\execution scenarios};
    \node[ormblock, below=of $(unit)!0.5!(integration)$] (cmake) {\path{tests/*/CMakeLists.txt}\\registration + gating};
    \node[ormblock, below=of cmake] (ctest) {\code{gtest\_discover\_tests}\\и \code{ctest}};

    \draw[ormline] (src) -- (unit);
    \draw[ormline] (src) -- (integration);
    \draw[ormline] (unit) -- (cmake);
    \draw[ormline] (integration) -- (cmake);
    \draw[ormline] (cmake) -- (ctest);
\end{tikzpicture}
\caption{Верификационен слой на хранилището и поток към изпълнимите тестове}
\label{fig:verification_appendix_map}
\end{figure}

Фигура~\ref{fig:verification_appendix_map} показва защо verification-ът не трябва да се възприема като единичен тестов пакет. Той е съставен от няколко нива: локални свойства на типовете, contract-но поведение на connector-ите, интеграционни сценарии за query execution и build-level механизъм, който решава кои тестове са налични в конкретната среда.

\subsection{Каталог на основните unit тестове}

Unit тестовете покриват две различни категории знания. Първата категория е свързана с логическия модел и с основните generic механизми на библиотеката. Втората категория е свързана с backend-специфичните connector-и и техните capability профили.

\begin{table}[H]
\centering
\small
\caption{Основни unit тестове за generic и core подсистемите}
\label{tab:appendix_unit_core}
\begin{tabularx}{\textwidth}{L{3.3cm}L{3.8cm}Y}
\toprule
\textbf{Файл} & \textbf{Подсистема} & \textbf{Какво валидира} \\
\midrule
\path{test_types.cpp} & базови типови alias-и & коректност на публичните alias-и, включително chrono-related типове \\
\path{test_property.cpp} & entity property layer & compile-time модел на скаларните полета и type trait поведението \\
\path{test_relationship.cpp} & entity relationship layer & infer логика за връзките и storage strategy поведението \\
\path{test_query.cpp} & query IR и fluent API & коректност на \code{select}, \code{insert}, \code{update}, \code{deleteq} и производните query структури \\
\path{test_thread_safety.cpp} & concurrency helpers & поведение на connection pool, thread-local достъп и transaction guard механизма \\
\path{test_wire_protocol.cpp} & experimental execution layer & ниско ниво на compile-time/wire-oriented механизми \\
\path{test_cpp26_reflection.cpp} & future language integration & early path за reflection-базирана автоматизация \\
\path{test_schema_migration.cpp} & migration layer & DDL generation, diff логика и state machine поведение \\
\bottomrule
\end{tabularx}
\end{table}

Таблица~\ref{tab:appendix_unit_core} е важна, защото показва, че unit слоят не се изчерпва с connector тестове. Той покрива и generic архитектурните елементи, върху които по-късно стъпват всички backend-и.

\begin{table}[H]
\centering
\small
\caption{Backend-oriented unit тестове и техният contract focus}
\label{tab:appendix_unit_connectors}
\begin{tabularx}{\textwidth}{L{3.4cm}L{3.3cm}Z}
\toprule
\textbf{Файл} & \textbf{Backend фокус} & \textbf{Тип доказателство} \\
\midrule
\path{test_postgresql_connector.cpp} & \code{PostgreSQLDB} & \code{is\_connector} compliance, capability asserts и SQL parameter rendering \\
\path{test_mysql_connector.cpp} & \code{MySQLDB} & joins/transactions capability профил и MySQL-specific rendering expectations \\
\path{test_mongodb_connector.cpp} & \code{MongoDB} & negative capability validation и BSON-like filter rendering \\
\path{test_redis_connector.cpp} & \code{RedisDB} & absence of joins/transactions/aggregation и key-value command materialization \\
\path{test_cassandra_connector.cpp} & \code{CassandraDB} & capability ограничения и CQL rendering корелати \\
\path{test_neo4j_connector.cpp} & \code{Neo4jDB} & transaction capability и graph-oriented query materialization \\
\bottomrule
\end{tabularx}
\end{table}

\subsection{Integration и live тестове}

Integration слоят е организиран около два принципа. Първият е да има backend-независим mock baseline, който може да валидира query execution contract-а без външна инфраструктура. Вторият е да има live тестове за реални backend-и, когато съответните targets и environment flags са налични.

\begin{table}[H]
\centering
\small
\caption{Каталог на integration и live тестовете}
\label{tab:appendix_integration_tests}
\begin{tabularx}{\textwidth}{L{3.2cm}L{3.0cm}Y}
\toprule
\textbf{Файл} & \textbf{Режим на изпълнение} & \textbf{Какво демонстрира} \\
\midrule
\path{test_mockdb.cpp} & винаги наличен integration baseline & query composition, parameter binding, prepared queries и CRUD сценарии без външен backend \\
\path{test_sqlite.cpp} & локален integration тест & реален релационен execution path с capability asserts и hydration сценарии \\
\path{test_postgresql_live.cpp} & условен live тест & реална връзка, table lifecycle и CRUD изпълнение върху PostgreSQL \\
\path{test_mysql_live.cpp} & условен live тест & реален MySQL driver path, insert/update/delete и select filtering \\
\path{test_mongodb_live.cpp} & условен live тест & изпълнение на document-oriented сценарии срещу жив backend \\
\path{test_redis_live.cpp} & условен live тест & key-value и hash command materialization върху Redis среда \\
\path{test_cassandra_live.cpp} & условен live тест & wide-column execution и backend-specific CQL сценарии \\
\path{test_neo4j_live.cpp} & условен live тест & graph query execution през контролирана Docker среда \\
\bottomrule
\end{tabularx}
\end{table}

Този каталог е важен, защото показва, че repository-backed доказателството е разслоено. Не всички backend-и имат еднакъв operational cost за live verification, но хранилището съдържа ясна стратегия как и кога подобни тестове се активират.

\subsection{Capability доказателства по backend семейства}

Една от най-съществените идеи на библиотеката е, че connector-ите не трябва да претендират за повече възможности, отколкото реално предоставят. Затова capability доказателствата са разпределени между unit tests и integration tests, а не само между документационни твърдения.

\begin{table}[H]
\centering
\small
\caption{Примери за capability доказателства в тестовия слой}
\label{tab:appendix_capability_evidence}
\begin{tabularx}{\textwidth}{L{2.6cm}L{2.6cm}Z}
\toprule
\textbf{Backend} & \textbf{Основен тестов файл} & \textbf{Какво се доказва} \\
\midrule
SQLite & \path{test_sqlite.cpp} & налични \code{supports\_joins} и \code{supports\_transactions} в реален integration контекст \\
PostgreSQL & \path{test_postgresql_connector.cpp} & налични joins, transactions и aggregation capability tags \\
MySQL & \path{test_mysql_connector.cpp} & налични joins и transactions, но не и aggregation capability \\
MongoDB & \path{test_mongodb_connector.cpp} & отсъствие на SQL-oriented capability tags и коректно filter rendering поведение \\
Redis & \path{test_redis_connector.cpp} & отсъствие на joins/transactions/aggregation и key-value materialization semantics \\
Neo4j & \path{test_neo4j_connector.cpp} & налични transactions и graph-specific execution focus \\
Cassandra & \path{test_cassandra_connector.cpp} & ограничен capability профил и wide-column contract поведение \\
\bottomrule
\end{tabularx}
\end{table}

\subsection{CMake registration като част от verification дизайна}

Build системата не е неутрален фон, а активен компонент на верификацията. Това се вижда ясно от \path{tests/unit/CMakeLists.txt} и \path{tests/integration/CMakeLists.txt}. В unit слоя всички core тестове са обединени в един изпълним target и се регистрират чрез \code{gtest\_discover\_tests}. В integration слоя част от тестовете са винаги налични, а live тестовете се активират само когато съответният target и feature flag са налице.

\begin{lstlisting}[caption={Извадка от unit test registration-а},label={lst:appendix_unit_cmake}]
add_executable(test_unit
    test_types.cpp
    test_property.cpp
    test_relationship.cpp
    test_query.cpp
    test_mysql_connector.cpp
    test_postgresql_connector.cpp
    test_mongodb_connector.cpp
    test_redis_connector.cpp
    test_cassandra_connector.cpp
    test_neo4j_connector.cpp
    test_thread_safety.cpp
    test_wire_protocol.cpp
    test_cpp26_reflection.cpp
    test_schema_migration.cpp
)

gtest_discover_tests(test_unit)
\end{lstlisting}

Listing~\ref{lst:appendix_unit_cmake} е кратък, но много показателен: той описва unit verification слоя като систематичен каталог, а не като спорадично добавени тестове.

\begin{lstlisting}[caption={Извадка от conditional live test registration-а},label={lst:appendix_integration_cmake}]
add_executable(test_integration
    test_mockdb.cpp
)

gtest_discover_tests(test_integration)

if(TARGET orm::postgresql_live AND ORM_POSTGRESQL_LIVE_ENABLED)
    add_executable(test_postgresql_live test_postgresql_live.cpp)
    add_test(NAME PostgreSQLLive COMMAND test_postgresql_live)
endif()

if(TARGET orm::mysql_live AND ORM_MYSQL_LIVE_ENABLED)
    add_executable(test_mysql_live test_mysql_live.cpp)
    add_test(NAME MySQLLive COMMAND test_mysql_live)
endif()
\end{lstlisting}

Listing~\ref{lst:appendix_integration_cmake} показва, че integration/live verification-ът е преднамерено conditional. Това е архитектурно разумно решение, защото не всеки backend е винаги наличен във всяка среда, но framework-ът за регистрация остава формално определен.

\subsection{Изводи от приложението}

Настоящото приложение показва, че верификацията на библиотеката има собствена вътрешна архитектура. Тя включва каталог на unit доказателства, baseline integration execution path, условни live сценарии и build-level механизъм за регистрация и активиране на тестовете. Това е важна част от изследователската стойност на разработката, защото прави възможно твърденията от основния текст да бъдат свързани с конкретни, проверими артефакти.

Следователно дипломната работа не просто описва library design, а документира система, чиито основни тези могат да бъдат проследени до source headers, CMake registration и изпълними тестове. Именно тази проследимост е критична за оценяването на \ORMName{} библиотеката като завършена инженерна и академична разработка.
